This thesis expands robotic manipulation beyond geometric reasoning by developing methods that model and plan over dynamic and non-prehensile interactions, enabling new capabilities governed by object inertial properties, impulsive contact, and motion-induced forces.

Across three parts---online kinodynamic search, amortized reusable motion structure, and learned conditional generative models---the work studies where knowledge of dynamic feasibility should reside in a manipulation system, what form it should take, and how that choice shapes the behaviors a robot can generate.
Each part develops a different mechanism for placing constraint knowledge: an online planner that evaluates dynamics through simulation, a precomputed graph that amortizes dynamic structure across queries, and a learned distribution that absorbs feasibility into the generation process itself.

This chapter first revisits the main contributions and then distills the broader lessons that cut across them.
Those lessons fall into three classes---constraints and system design, multimodal capabilities beyond geometric skills, and the scope and limits of statistical learning---each paired with the open questions it raises for the next generation of dynamics-aware manipulation systems.
